% from ega0-1-fr.tex
\oldpage[0\textsubscript{I}]{27}

de $\mathfrak{B}$ contenus dans un $U_\alpha$ au moins~; il est clair que
$\mathfrak{B}'$ est encore une base de la topologie de $X$, donc
$\mathcal{F}'(U)$ (resp. $\mathcal{F}''(U_\alpha)$) est limite projective des
$\mathcal{F}(V)$ pour $V\in\mathfrak{B}'$ et $V\subset U$ (resp.
$V\subset U_\alpha$)~; l'axiome (F) se vérifie alors aussitôt en vertu de la
définition de la limite projective.

Lorsque (F\textsubscript{0}) est vérifié, nous dirons par abus de langage que
le préfaisceau $\mathcal{F}$ sur la base $\mathfrak{B}$ est un \emph{faisceau}.
\end{env}

\begin{env}[3.2.3]
\label{0.3.2.3-fr}
Soient $\mathcal{F}$, $\mathcal{G}$ deux préfaisceaux sur la base
$\mathfrak{B}$, à valeurs dans $\mathbf{K}$~; on définit un \emph{morphisme}
$u:\mathcal{F}\to\mathcal{G}$ comme une famille
$(u_V)_{V\in\mathfrak{B}}$ de morphismes
$u_V:\mathcal{F}(V)\to\mathcal{G}(V)$ satisfaisant aux conditions de
compatibilité usuelles avec les morphismes de restriction $\rho_V^W$. Avec
les notations de (3.2.1), on en déduit un morphisme
$u':\mathcal{F}'\to\mathcal{G}'$ des préfaisceaux (ordinaires) correspondants
en prenant pour $u_U'$ la limite projective des $u_V$ pour
$V\in\mathfrak{B}$ et $V\subset U$~; la vérification des conditions de
compatibilité avec $\rho'_U{}^{U'}$ découle du caractère fonctoriel de la
limite projective.
\end{env}

\begin{env}[3.2.4]
\label{0.3.2.4-fr}
Si la catégorie $\mathbf{K}$ admet des limites inductives, et si $\mathcal{F}$
est un préfaisceau sur la base $\mathfrak{B}$, à valeurs dans $\mathbf{K}$,
pour tout $x\in X$ les voisinages de $x$ appartenant à $\mathfrak{B}$ forment
un ensemble cofinal (pour $\supset$) dans l'ensemble des voisinages de $x$,
donc, si $\mathcal{F}'$ est le préfaisceau (ordinaire) correspondant à
$\mathcal{F}$, la fibre $\mathcal{F}_x'$ est égale à
\[
  \varinjlim_{\mathfrak{B}}\mathcal{F}(V)
\]
selon l'ensemble des $V\in\mathfrak{B}$ contenant $x$.
Si $u:\mathcal{F}\to\mathcal{G}$ est un morphisme de préfaisceaux sur
$\mathfrak{B}$ à valeurs dans $\mathbf{K}$,
$u':\mathcal{F}'\to\mathcal{G}'$ le morphisme correspondant de préfaisceaux
ordinaires, $u_x'$ est de même la limite inductive des morphismes
$u_V:\mathcal{F}(V)\to\mathcal{G}(V)$ pour $V\in\mathfrak{B}$, $x\in V$.
\end{env}

\begin{env}[3.2.5]
\label{0.3.2.5-fr}
Revenons aux conditions générales de (3.2.1). Si $\mathcal{F}$ est un
\emph{faisceau} ordinaire à valeurs dans $\mathbf{K}$, $\mathcal{F}_1$ le
faisceau \emph{sur $\mathfrak{B}$} obtenu par restriction de $\mathcal{F}$ à
$\mathfrak{B}$, le faisceau ordinaire $\mathcal{F}_1'$ obtenu à partir de
$\mathcal{F}_1$ par le procédé de (3.2.1) est canoniquement isomorphe à
$\mathcal{F}$, en vertu de la condition (F) et des propriétés d'unicité de la
limite projective. On identifiera d'ordinaire $\mathcal{F}$ et
$\mathcal{F}_1'$.

Si $\mathcal{G}$ est un second faisceau (ordinaire) sur $X$ à valeurs dans
$\mathbf{K}$, et $u:\mathcal{F}\to\mathcal{G}$ un morphisme, la remarque
précédente montre que la donnée des
$u_V:\mathcal{F}(V)\to\mathcal{G}(V)$ \emph{pour les $V\in\mathfrak{B}$
seulement} détermine complètement $u$~; inversement, il suffit, les $u_V$
étant donnés pour $V\in\mathfrak{B}$, de vérifier le diagramme de commutativité
avec les morphismes de restriction $\rho_V^W$ pour $V\in\mathfrak{B}$,
$W\in\mathfrak{B}$ et $V\subset W$, pour qu'il existe un morphisme $u'$ et un
seul de $\mathcal{F}$ dans $\mathcal{G}$ tel que $u_V'=u_V$ pour tout
$V\in\mathfrak{B}$ (3.2.3).
\end{env}

\begin{env}[3.2.6]
\label{0.3.2.6-fr}
Supposons toujours que $\mathbf{K}$ admette des limites projectives. Alors la
catégorie des \emph{faisceaux sur $X$ à valeurs dans $\mathbf{K}$} admet aussi
des \emph{limites projectives}~; si $(\mathcal{F}_\lambda)$ est un système
projectif de faisceaux sur $X$ à valeurs dans $\mathbf{K}$, les
\[
  \mathcal{F}(U)=\varprojlim_\lambda\mathcal{F}_\lambda(U)
\]
définissent en effet un préfaisceau à valeurs dans $\mathbf{K}$, et la
vérification de l'axiome (F) résulte de la transitivité des limites
projectives~; le fait que $\mathcal{F}$ est alors limite projective des
$\mathcal{F}_\lambda$ est immédiat.

Lorsque $\mathbf{K}$ est la catégorie des ensembles, pour tout système
projectif $(\mathcal{H}_\lambda)$ tel

